\documentclass[11pt]{article}
\usepackage[a4paper,margin=1in]{geometry}
\usepackage{amsmath,amsthm,amssymb}
\usepackage{mathtools}

\newtheorem{theorem}{Theorem}[section]
\newtheorem{proposition}[theorem]{Proposition}
\newtheorem{lemma}[theorem]{Lemma}
\newtheorem{corollary}[theorem]{Corollary}
\theoremstyle{definition}
\newtheorem{definition}[theorem]{Definition}
\newtheorem{remark}[theorem]{Remark}
\newtheorem{hypothesis}[theorem]{Hypothesis}

\newcommand{\R}{\mathbb{R}}
\newcommand{\Fr}{\mathcal{A}}
\newcommand{\Om}{\Omega}
\newcommand{\eps}{\varepsilon}
\newcommand{\dt}[1]{\frac{d#1}{d\rho}}
\newcommand{\rstar}{\rho_{*}}

\title{Exact Deficit Identity on Reduced-Boundary Level Sets\\
       (Route~$\delta$, Plan~2, Building Block \textsc{LevelSetIdentity})}
\author{}
\date{}

\begin{document}
\maketitle

\begin{abstract}
We prove that for every open set $\Om\subset\R^{n}$ of finite measure, with
$u=u_{\Om}\in H^{1}_{0}(\Om)$ the variational torsion function and
$E_{t}:=\{u>t\}$, $m(t):=|E_{t}|$, the Saint--Venant deficit admits the
exact representation
\[
   \frac{2}{|\Om|}\bigl(E(\Om)-E(B)\bigr)
   =\frac{1}{|\Om|}\int_{0}^{\|u\|_{\infty}}
     \frac{D_{H}(t)+D_{I}(t)}{n^{2}\omega_{n}^{2/n}\,m(t)^{1-2/n}}\,dt,
\]
where $B$ is the ball with $|B|=|\Om|$, $D_{H}(t)$ is the Cauchy--Schwarz
defect of $|\nabla u|$ on the reduced boundary $\partial^{*}E_{t}$, and
$D_{I}(t)$ is the squared-perimeter isoperimetric defect of $E_{t}$. The
proof rests on four ingredients which all hold on reduced boundaries for
a.e.\ level: (i) the BV coarea identity for $-m'(t)$, (ii) a weak
divergence-flux identity proved by Lipschitz truncation, (iii) the
distributional second derivative of the rearranged primitive $I(s)=\int_{0}^{s}
u^{*}$, in the corrected form $I''=-1/\alpha$ with
$\alpha=\int_{\partial^{*}E_{t}}|\nabla u|^{-1}$, and (iv) the endpoint
convexity gap together with the P\'olya--Szeg\H{o} identification of the
ball energy. No graph regularity, no BDV selection, no smoothness of
$\partial\Om$, and no boundedness or connectedness of $\Om$ are used. The
identity is the foundational input for the downstream Plan~2 blocks
\cite{MetricFramework,SpaceTimeIdentity,WeightedMetricTrace,
BoundaryLayerTransfer,GlobalAssembly}.
\end{abstract}

\section{Notation and standing hypotheses}\label{sec:notation}

Throughout $n\ge 2$. We fix the following objects and use them without
further comment.

\begin{itemize}
\item $\Om\subset\R^{n}$ open with $|\Om|<\infty$. No connectedness,
smoothness, or boundedness is assumed.

\item $u=u_{\Om}\in H^{1}_{0}(\Om)$ is the variational torsion function,
i.e.\ the unique minimizer of $J(v)=\tfrac12\int|\nabla v|^{2}-\int v$
over $H^{1}_{0}(\Om)$; equivalently the weak solution of
$-\Delta u=1$ in $\Om$, $u=0$ on $\partial\Om$ (in the $H^{1}_{0}$ trace
sense). Testing against $u^{-}$ yields $u\ge 0$ a.e.\ by the weak
maximum principle.

\item $E(\Om):=-\tfrac12\int_{\Om}u\,dx=-\tfrac12 T(\Om)$, where
$T(\Om)$ is the torsional rigidity.

\item $B\subset\R^{n}$ is the open ball with $|B|=|\Om|$, radius
$R:=(|\Om|/\omega_{n})^{1/n}$, $\omega_{n}=|B_{1}|$. The function
$u_{B}$ is the torsion function of $B$.

\item For $t\ge 0$, $E_{t}:=\{u>t\}$ (the precise representative as in
\cite[\S 15]{Maggi2012}), $m(t):=|E_{t}|$. The function $m$ is
right-continuous, nonincreasing, with $m(0^{+})\le|\Om|$ and
$m(\|u\|_{\infty})=0$.

\item $u^{*}\colon[0,|\Om|]\to[0,\|u\|_{\infty}]$ is the decreasing
rearrangement, $u^{*}(s):=\inf\{t\ge 0:m(t)\le s\}$, in the convention of
\cite{Kawohl1985,Talenti1976}.

\item $\partial^{*}E_{t}$ denotes the De~Giorgi reduced boundary, with
$\mathcal H^{n-1}(\partial E_{t}\setminus\partial^{*}E_{t})=0$ for a.e.\
$t$ \cite{DeGiorgi1958}.

\item We set
\[
   \alpha(t):=\int_{\partial^{*}E_{t}}|\nabla u|^{-1}\,d\mathcal H^{n-1},
   \qquad
   \beta(t):=\int_{\partial^{*}E_{t}}|\nabla u|\,d\mathcal H^{n-1},
\]
defined for a.e.\ $t$ (see Lemmas~\ref{lem:coarea} and~\ref{lem:flux}), and
$c_{n}:=n^{2}\omega_{n}^{2/n}$.
\end{itemize}

\paragraph{Talenti $L^{\infty}$ bound.} By the Talenti pointwise
comparison \cite{Talenti1976},
\begin{equation}\label{eq:talenti}
   u^{*}_{\Om}(s)\le u^{*}_{B}(s)
   =\frac{R^{2}-(s/\omega_{n})^{2/n}}{2n},
   \qquad s\in[0,|\Om|],
\end{equation}
so $\|u\|_{L^{\infty}(\Om)}\le R^{2}/(2n)<\infty$. This is the only
reason boundedness of $\Om$ is not needed at the analytic level.

\paragraph{Coarea applicability.} Since $u\in H^{1}_{0}(\Om)\subset
BV_{\rm loc}(\R^{n})$ (extended by $0$) and $u\in L^{\infty}$, the BV
coarea formula \cite[Thm.~13.1]{Maggi2012}, \cite[Thm.~3.40]{AFP2000}
gives, for every Borel $\varphi\ge 0$,
\begin{equation}\label{eq:coarea}
   \int_{\Om}|\nabla u|\,\varphi\,dx
   =\int_{-\infty}^{\infty}\Bigl(\int_{\partial^{*}E_{t}}\varphi\,d\mathcal
   H^{n-1}\Bigr)\,dt.
\end{equation}
The set
\[
   \mathcal R:=\bigl\{t\in(0,\|u\|_{\infty}):
   P(E_{t})<\infty,\ \partial^{*}E_{t}\subset\Om\bigr\}
\]
has full Lebesgue measure in $(0,\|u\|_{\infty})$; Sobolev--Sard
\cite[App.~A, Thm.~A.1]{FigalliMaggi2011} additionally gives
$\mathcal H^{n-1}(\{u=t\}\cap\{|\nabla u|=0\})=0$ for a.e.\ $t$, but is
not used below since all integrals are on reduced boundaries.
Statements ``for a.e.\ $t$'' below refer to a.e.\ $t\in\mathcal R$.

\section{Coarea identity for $-m'(t)$}\label{sec:coarea}

\begin{lemma}[Coarea form of $-m'$]\label{lem:coarea}
For a.e.\ $t\in(0,\|u\|_{\infty})$,
\begin{equation}\label{eq:L1}
   -m'(t)
   =\int_{\partial^{*}E_{t}}\frac{1}{|\nabla u|}\,d\mathcal H^{n-1}
   =\alpha(t).
\end{equation}
In particular $m$ is absolutely continuous on $[\eps,\|u\|_{\infty}]$ for
every $\eps>0$, and $\alpha\in L^{1}_{\rm loc}((0,\|u\|_{\infty}))$.
\end{lemma}

\begin{proof}
Fix $0<\alpha_{0}<\beta_{0}<\|u\|_{\infty}$ and apply~\eqref{eq:coarea}
to the Borel function $\varphi=|\nabla u|^{-1}
\mathbf 1_{\{|\nabla u|>0\}}\mathbf 1_{[\alpha_{0},\beta_{0}]}(u)$. By
\cite[Thm.~13.1]{Maggi2012},
\[
   \int_{\Om}\mathbf 1_{\{\alpha_{0}<u<\beta_{0}\}}\,dx
   =\int_{\alpha_{0}}^{\beta_{0}}
     \Bigl(\int_{\partial^{*}E_{t}}|\nabla u|^{-1}\,d\mathcal H^{n-1}
     \Bigr)dt,
\]
because the integrand on the left equals
$|\nabla u|\cdot|\nabla u|^{-1}\mathbf 1_{\{|\nabla u|>0\}}
\mathbf 1_{\{\alpha_{0}<u<\beta_{0}\}}$ and the contribution of
$\{|\nabla u|=0\}$ to the BV coarea slicing vanishes \cite[3.40(b)]{AFP2000}.
The left-hand side equals $m(\alpha_{0})-m(\beta_{0})$ outside the at
most countably many plateau levels of $u$ (levels $t$ with $|\{u=t\}|>0$;
each such plateau contributes a jump of $m$, and $m$ is monotone, so
these are at most countably many). Differentiating in $\beta_{0}$ at
a.e.\ $\beta_{0}$ yields~\eqref{eq:L1}.
\end{proof}

\begin{remark}\label{rem:plateaus}
At the countable set of plateau levels~\eqref{eq:L1} is not asserted; the
Lebesgue integral in the main theorem absorbs them trivially.
\end{remark}

\section{Weak divergence-flux identity}\label{sec:flux}

\begin{lemma}[Flux $=$ volume]\label{lem:flux}
For a.e.\ $t\in(0,\|u\|_{\infty})$, $E_{t}$ has finite perimeter and
\begin{equation}\label{eq:L2}
   \beta(t)=
   \int_{\partial^{*}E_{t}}|\nabla u|\,d\mathcal H^{n-1}
   =m(t).
\end{equation}
\end{lemma}

\begin{proof}
Fix $t\in\mathcal R$ with $|E_{t}|<\infty$, $P(E_{t})<\infty$, and $u$
not having a plateau at $t$ (all these hold for a.e.\ $t$).

\smallskip\noindent
\emph{Step 1 (Lipschitz truncation).}
For each $\eps>0$ set
\[
   \phi_{\eps}(s):=\min\bigl(1,\max(0,s/\eps)\bigr),
   \qquad
   \eta_{\eps}(x):=\phi_{\eps}\bigl(u(x)-t\bigr).
\]
Since $\phi_{\eps}$ is Lipschitz with $\phi_{\eps}(0)=0$ and
$u\in H^{1}_{0}(\Om)$, one has $\eta_{\eps}\in H^{1}_{0}(\Om)$. Testing
the weak equation $-\Delta u=1$ in $H^{1}_{0}(\Om)$ against $\eta_{\eps}$
gives
\begin{equation}\label{eq:weakform}
   \int_{\Om}\nabla u\cdot\nabla\eta_{\eps}\,dx
   =\int_{\Om}\eta_{\eps}\,dx.
\end{equation}
The chain rule for Sobolev compositions \cite[Thm.~3.99]{AFP2000} gives
$\nabla\eta_{\eps}=\eps^{-1}\nabla u\,\mathbf 1_{\{t<u<t+\eps\}}$, so
\[
   \int_{\Om}\nabla u\cdot\nabla\eta_{\eps}\,dx
   =\frac{1}{\eps}\int_{\{t<u<t+\eps\}}|\nabla u|^{2}\,dx.
\]

\smallskip\noindent
\emph{Step 2 (Coarea on the right of~\eqref{eq:weakform}).}
Apply~\eqref{eq:coarea} with
$\varphi=|\nabla u|\,\mathbf 1_{(t,t+\eps)}\circ u$:
\[
   \int_{\{t<u<t+\eps\}}|\nabla u|^{2}\,dx
   =\int_{t}^{t+\eps}\Bigl(\int_{\partial^{*}E_{\tau}}
      |\nabla u|\,d\mathcal H^{n-1}\Bigr)d\tau
   =\int_{t}^{t+\eps}\beta(\tau)\,d\tau.
\]

\smallskip\noindent
\emph{Step 3 (Limit on the left of~\eqref{eq:weakform}).}
$\eta_{\eps}\to\mathbf 1_{E_{t}}$ pointwise as $\eps\downarrow 0$
(away from the level set $\{u=t\}$, which has Lebesgue measure zero
since $t$ is not a plateau) and $0\le\eta_{\eps}\le 1$, so by dominated
convergence $\int_{\Om}\eta_{\eps}\to|E_{t}|=m(t)$.

\smallskip\noindent
\emph{Step 4 (Limit on the right).}
The map $\tau\mapsto\beta(\tau)$ belongs to $L^{1}_{\rm loc}((0,\|u\|_{\infty}))$:
by~\eqref{eq:coarea} with $\varphi=|\nabla u|\,\mathbf 1_{(0,\|u\|_{\infty})}\circ u$,
\[
   \int_{0}^{\|u\|_{\infty}}\beta(\tau)\,d\tau
   =\int_{\Om}|\nabla u|^{2}\,dx
   =\int_{\Om}u\,dx<\infty,
\]
the last identity being the Dirichlet identity obtained by testing
$-\Delta u=1$ against $u\in H^{1}_{0}(\Om)$. At every Lebesgue point of
$\beta$, hence for a.e.\ $t$,
\[
   \frac{1}{\eps}\int_{t}^{t+\eps}\beta(\tau)\,d\tau
   \xrightarrow[\eps\downarrow 0]{}\beta(t).
\]

\smallskip\noindent
\emph{Step 5.} Combining Steps 1--4 with~\eqref{eq:weakform} yields
$\beta(t)=m(t)$ for a.e.\ $t$.
\end{proof}

\begin{remark}[Formal divergence-theorem reading]\label{rem:divergence}
The argument of Lemma~\ref{lem:flux} is the rigorous BV/Lipschitz-truncation
replacement for the formal divergence theorem on $E_{t}$:
$\nabla u\cdot\nu_{E_{t}}=-|\nabla u|$ on $\partial^{*}E_{t}$ (since
$u>t$ inside), so $-\Delta u=1$ would give
\[
   m(t)=\int_{E_{t}}1\,dx=-\int_{E_{t}}\Delta u\,dx
   =-\int_{\partial^{*}E_{t}}\nabla u\cdot\nu_{E_{t}}\,d\mathcal H^{n-1}
   =\beta(t).
\]
Without a boundary trace of $\nabla u$ on $\partial^{*}E_{t}$, the
Lipschitz truncation above is what makes this rigorous on arbitrary
finite-perimeter level sets \cite[Ch.~17]{Maggi2012}.
\end{remark}

\begin{remark}[Integrated Dirichlet identity]
Integrating~\eqref{eq:L2} against $dt$ and using~\eqref{eq:coarea}
recovers the Dirichlet identity
$\int_{0}^{\|u\|_{\infty}}\beta(t)\,dt=\int_{\Om}|\nabla u|^{2}
=\int_{\Om}u=2|E(\Om)|$ used above. Lemma~\ref{lem:flux} is its
pointwise level-wise version.
\end{remark}

\section{Distributional derivatives of the rearranged primitive}
\label{sec:Iprimitive}

Define
\begin{equation}\label{eq:Idef}
   I(s):=\int_{\{u>u^{*}(s)\}}u\,dx,\qquad s\in[0,|\Om|].
\end{equation}
The following three lemmas identify $I$ with a one-dimensional primitive
of $u^{*}$ and compute its second distributional derivative in the form
used by the deficit identity.

\begin{lemma}[Layer-cake form of $I$]\label{lem:I3a}
For every $s\in[0,|\Om|]$,
\begin{equation}\label{eq:L3a}
   I(s)=\int_{0}^{s}u^{*}(\sigma)\,d\sigma.
\end{equation}
\end{lemma}

\begin{proof}
By equimeasurability of $u$ and $u^{*}$ on $(0,|\Om|)$
\cite[\S 2]{Kawohl1985}, \cite{Talenti1976},
$\int_{\Om}g(u)\,dx=\int_{0}^{|\Om|}g(u^{*}(\sigma))\,d\sigma$ for every
Borel $g\ge 0$. Choose $g(t)=t\,\mathbf 1_{\{t>u^{*}(s)\}}$; since
$\{\sigma\in(0,|\Om|):u^{*}(\sigma)>u^{*}(s)\}=(0,s)$ up to null sets by
the very definition of $u^{*}$, we obtain
\[
   I(s)=\int_{\Om}g(u)\,dx
       =\int_{0}^{|\Om|}g(u^{*}(\sigma))\,d\sigma
       =\int_{0}^{s}u^{*}(\sigma)\,d\sigma.
\]
\end{proof}

\begin{lemma}[First derivative]\label{lem:I3b}
$I$ is absolutely continuous on $[0,|\Om|]$ and $I'(s)=u^{*}(s)$ for
a.e.\ $s\in(0,|\Om|)$.
\end{lemma}

\begin{proof}
By~\eqref{eq:talenti} $u^{*}\in L^{\infty}(0,|\Om|)$; in particular
$u^{*}$ is bounded and Borel, hence in $L^{1}(0,|\Om|)$ (since
$|\Om|<\infty$). By
Lemma~\ref{lem:I3a}, $I$ is the indefinite integral of $u^{*}$, so the
Lebesgue differentiation theorem gives the claim.
\end{proof}

\begin{lemma}[Second derivative --- corrected form]\label{lem:I3c}
The rearrangement $u^{*}$ is absolutely continuous on
$[s_{0},|\Om|-\eps]$ for every $0<s_{0}<|\Om|-\eps<|\Om|$, and

Two intermediate steps justify the displayed chain below: (a) the
identity $I''=(u^{*})'$ follows from Lemma~\ref{lem:I3b}
($I'=u^{*}$) combined with the absolute continuity of $u^{*}$ on
closed subintervals of $(0,|\Om|)$ (so that the distributional
derivative of $I'$ coincides with the a.e.\ classical derivative of
$u^{*}$); and (b) $\alpha(t)>0$ for a.e.\ $t$, since $\alpha=-m'$
by Lemma~\ref{lem:coarea} and $m$ is strictly decreasing a.e.\ on
its non-plateau range, so the reciprocal is well defined a.e.
\begin{equation}\label{eq:L3c}
   I''(s)=(u^{*})'(s)
   =\frac{1}{m'(u^{*}(s))}
   =-\frac{1}{\alpha(u^{*}(s))}
   \qquad\text{for a.e.\ }s\in(0,|\Om|).
\end{equation}
\end{lemma}

\begin{proof}
$m$ is right-continuous, nonincreasing, and (Lemma~\ref{lem:coarea})
absolutely continuous on $[\eps,\|u\|_{\infty}]$ for every $\eps>0$,
with $m'(t)=-\alpha(t)$ a.e. The plateau levels of $u$ (countable, by
monotonicity of $m$) are the levels where $m$ jumps; outside them, $m$
is strictly decreasing on its domain of strict monotonicity, hence its
right-continuous generalized inverse $u^{*}$ is also strictly
decreasing and absolutely continuous on closed subintervals of
$(0,|\Om|)$. Differentiating $m(u^{*}(s))=s$ at a.e.\ $s$ via the
inverse-function theorem for monotone AC functions \cite[Thm.~3.99]{AFP2000},
\cite{BBC1975} gives~\eqref{eq:L3c}.
\end{proof}

\begin{remark}[Notational fix vs.\ the source note]\label{rem:fix}
The companion note `level-set-deficit-identity.md' (audit-trail
markdown note, not a mathematical reference) wrote the second
derivative as $I''(s)=-1/a(u^{*}(s))$ with $a(t):=\int_{\partial^{*}E_{t}}
|\nabla u|$, i.e.\ with $\beta$ in place of $\alpha$. By
Lemma~\ref{lem:flux} we have $\beta(t)=m(t)$, so this would give
$I''=-1/m$ --- correct only when $|\nabla u|$ is constant on
$\partial^{*}E_{t}$, i.e.\ on the ball, where Cauchy--Schwarz is sharp.
The general identity is~\eqref{eq:L3c}; the discrepancy between
$\alpha$ and $\beta$ is precisely the Cauchy--Schwarz defect $D_{H}$
introduced below, and replacing $\alpha$ by $\beta$ in $I''$ would
silently delete $D_{H}$ from the deficit identity. We use
the $\alpha$ form throughout.
\end{remark}

\paragraph{Cauchy--Schwarz.} The two integrals
$\alpha(t)=\int_{\partial^{*}E_{t}}|\nabla u|^{-1}\,d\mathcal H^{n-1}$
and $\beta(t)=\int_{\partial^{*}E_{t}}|\nabla u|\,d\mathcal H^{n-1}$ are
linked by Cauchy--Schwarz on $\partial^{*}E_{t}$ applied to the pair
$(|\nabla u|^{1/2},|\nabla u|^{-1/2})$,
\begin{equation}\label{eq:CS}
   P(E_{t})^{2}=\Bigl(\int_{\partial^{*}E_{t}}1\,d\mathcal H^{n-1}\Bigr)^{2}
   \le\alpha(t)\beta(t),
\end{equation}
with equality iff $|\nabla u|$ is $\mathcal H^{n-1}$-a.e.\ constant on
$\partial^{*}E_{t}$.

\section{Endpoint convexity gap and the deficit identity}
\label{sec:main}

Following the Nicola--Tilli convexity-gap idea adapted to torsion,
introduce
\begin{equation}\label{eq:Gdef}
   G(s):=I(s)+\frac{s^{1+2/n}}{(4+2n)\,\omega_{n}^{2/n}},
   \qquad s\in[0,|\Om|].
\end{equation}
The second derivative of the polynomial corrector is computed directly:
$\frac{d^{2}}{ds^{2}}\frac{s^{1+2/n}}{(4+2n)\omega_{n}^{2/n}}
=\frac{1}{c_{n}s^{1-2/n}}$, with $c_{n}=n^{2}\omega_{n}^{2/n}$.

\begin{lemma}[Pointwise convexity formula]\label{lem:Gpp}
For a.e.\ $s\in(0,|\Om|)$, with $t=u^{*}(s)$,
\begin{equation}\label{eq:Gpp}
   G''(s)
   =\frac{1}{c_{n}\,s^{1-2/n}}-\frac{1}{\alpha(t)}.
\end{equation}
\end{lemma}

\begin{proof}
Lemmas~\ref{lem:I3b}, \ref{lem:I3c} give $G''=I''+\frac{1}{c_{n}s^{1-2/n}}
=-1/\alpha(t)+\frac{1}{c_{n}s^{1-2/n}}$, which is~\eqref{eq:Gpp}.
\end{proof}

\begin{definition}[Two pointwise defects]\label{def:DHDI}
For a.e.\ $t\in(0,\|u\|_{\infty})$ set
\begin{align}
   D_{H}(t)&:=\alpha(t)\,\beta(t)-P(E_{t})^{2}
            =\Bigl(\!\int_{\partial^{*}E_{t}}\!\!\tfrac{1}{|\nabla u|}\Bigr)
              \Bigl(\!\int_{\partial^{*}E_{t}}\!\!|\nabla u|\Bigr)-P(E_{t})^{2},
              \label{eq:DH-def}\\[1mm]
   D_{I}(t)&:=P(E_{t})^{2}-c_{n}\,m(t)^{2-2/n}.
              \label{eq:DI-def}
\end{align}
By~\eqref{eq:CS}, $D_{H}(t)\ge 0$, with equality iff $|\nabla u|$ is
constant on $\partial^{*}E_{t}$. By the De~Giorgi isoperimetric
inequality \cite{DeGiorgi1958}, $P(E_{t})\ge n\omega_{n}^{1/n}|E_{t}|^{1-1/n}$
with equality iff $E_{t}$ is (equivalent to) a ball, hence
$D_{I}(t)\ge 0$ with equality iff $E_{t}$ is a ball.
\end{definition}

\begin{lemma}[Pointwise-to-integrated convexity identity]\label{lem:cov}
With $L:=|\Om|$,
\begin{equation}\label{eq:cov}
   \int_{0}^{L}s\,G''(s)\,ds
   =\int_{0}^{\|u\|_{\infty}}\frac{D_{H}(t)+D_{I}(t)}{c_{n}\,m(t)^{1-2/n}}\,dt.
\end{equation}
\end{lemma}

\begin{proof}
From~\eqref{eq:Gpp},
\[
   \int_{0}^{L}s\,G''(s)\,ds
   =\int_{0}^{L}s\Bigl(\frac{1}{c_{n}s^{1-2/n}}-\frac{1}{\alpha(u^{*}(s))}\Bigr)ds.
\]
Change variable $s=m(t)$, $ds=m'(t)\,dt=-\alpha(t)\,dt$ on the
strictly-decreasing branch of $m$ (the at most countable plateau levels
of $u$ are an $\mathcal L^{1}$-null set in $t$, contributing zero):
\[
   \int_{0}^{L}s\,G''(s)\,ds
   =\int_{0}^{\|u\|_{\infty}}\!m(t)\Bigl(\frac{1}{c_{n}m(t)^{1-2/n}}-
      \frac{1}{\alpha(t)}\Bigr)\alpha(t)\,dt
   =\int_{0}^{\|u\|_{\infty}}\Bigl(
      \frac{m(t)\alpha(t)}{c_{n}m(t)^{1-2/n}}-m(t)\Bigr)dt.
\]
Use $m(t)=\beta(t)$ (Lemma~\ref{lem:flux}) and add and subtract
$P(E_{t})^{2}/(c_{n}m(t)^{1-2/n})$ in the integrand:
\[
   \frac{m(t)\alpha(t)}{c_{n}m(t)^{1-2/n}}-m(t)
   =\frac{\alpha(t)\beta(t)-c_{n}m(t)^{2-2/n}}{c_{n}m(t)^{1-2/n}}
   =\frac{D_{H}(t)+D_{I}(t)}{c_{n}m(t)^{1-2/n}}.
\]
Substituting yields~\eqref{eq:cov}.
\end{proof}

\paragraph{Endpoint identification.} With $L:=|\Om|$ define
\begin{equation}\label{eq:Gammadef}
   \Gamma(\Om):=G'(L)-\frac{G(L)}{L}.
\end{equation}

\begin{lemma}[Endpoint convexity-gap formula]\label{lem:gap}
$\Gamma(\Om)=\frac{1}{L}\int_{0}^{L}s\,G''(s)\,ds$.
\end{lemma}

\begin{proof}
$G\in W^{2,1}_{\rm loc}((0,L])$ and $G(0)=I(0)=0$. Integration by parts
on $(\eps,L)$:
\[
   \int_{\eps}^{L}s\,G''(s)\,ds
   =\bigl[s\,G'(s)\bigr]_{\eps}^{L}-\int_{\eps}^{L}G'(s)\,ds
   =L\,G'(L)-\eps\,G'(\eps)-G(L)+G(\eps).
\]
By Lemma~\ref{lem:I3b}, $G'(s)=u^{*}(s)+\frac{s^{2/n}}{2n\omega_{n}^{2/n}}$
(using $\frac{1+2/n}{4+2n}=\frac{1}{2n}$),
which is bounded; hence $\eps G'(\eps)\to 0$ and $G(\eps)\to G(0)=0$ as
$\eps\downarrow 0$. Dividing by $L$ gives the claim.
\end{proof}

\begin{lemma}[Endpoint value]\label{lem:endpoint}
\begin{equation}\label{eq:endpoint}
   \Gamma(\Om)=\frac{2}{|\Om|}\bigl(E(\Om)-E(B)\bigr).
\end{equation}
\end{lemma}

\begin{proof}
At $s=L$, $u^{*}(L)=0$. By Lemma~\ref{lem:I3a},
$I(L)=\int_{0}^{L}u^{*}=\int_{\Om}u\,dx$, so
\[
   \frac{G(L)}{L}=\frac{1}{L}\int_{\Om}u\,dx
                  +\frac{L^{2/n}}{(4+2n)\omega_{n}^{2/n}}.
\]
By Lemma~\ref{lem:I3b}, $I'(L)=u^{*}(L)=0$, so $G'(L)=
\frac{L^{2/n}}{2n\omega_{n}^{2/n}}$. Hence
\[
   \Gamma(\Om)
   =\frac{L^{2/n}}{2n\omega_{n}^{2/n}}
    -\frac{L^{2/n}}{(4+2n)\omega_{n}^{2/n}}
    -\frac{1}{L}\int_{\Om}u\,dx
   =\frac{L^{2/n}}{2n(n+2)\omega_{n}^{2/n}}-\frac{1}{L}\int_{\Om}u\,dx.
\]
(Note $4+2n=2(n+2)$, reconciling the constant $1/(4+2n)$ in the
corrector~\eqref{eq:Gdef} with $1/(2n(n+2))$ here.)
For the ball $B=B_{R}$ with $|B_{R}|=L$, the torsion function is
$u_{B}(x)=(R^{2}-|x|^{2})/(2n)$ (so $-\Delta u_{B}=1$, $u_{B}=0$ on
$\partial B_{R}$). Take $|B|=\omega_{n}$ (so $R=1$, $L=\omega_{n}$,
$L^{2/n}=\omega_{n}^{2/n}$); direct integration in polar coordinates
yields
\[
   \int_{B_{1}}u_{B}\,dx
   =\frac{1}{2n}\int_{B_{1}}(1-|x|^{2})\,dx
   =\frac{1}{2n}\cdot n\omega_{n}\int_{0}^{1}(1-r^{2})r^{n-1}\,dr
   =\frac{\omega_{n}}{2n(n+2)}.
\]
Substituting into $\Gamma$ (still with $L=\omega_{n}$),
$\frac{1}{L}\int_{B}u_{B}=\frac{1}{2n(n+2)}$ and
$\frac{L^{2/n}}{2n(n+2)\omega_{n}^{2/n}}=\frac{1}{2n(n+2)}$, hence
$\Gamma(B)=0$. By scaling, the same conclusion holds for arbitrary
$L>0$, giving
$\frac{L^{2/n}}{2n(n+2)\omega_{n}^{2/n}}=\frac{1}{L}\int_{B}u_{B}\,dx$.
Subtracting,
\[
   \Gamma(\Om)
   =\frac{1}{L}\Bigl(\int_{B}u_{B}\,dx-\int_{\Om}u\,dx\Bigr)
   =\frac{2}{L}\bigl(E(\Om)-E(B)\bigr),
\]
using $E(\Om)=-\tfrac12\int_{\Om}u$. The identification of $\Gamma(B)$
with the radial computation is the classical P\'olya--Szeg\H{o} content
\cite{PolyaSzego1951}, \cite[\S II.5]{Bandle1980}.
\end{proof}

\begin{theorem}[Exact deficit identity on reduced-boundary level sets]
\label{thm:main}
Let $\Om\subset\R^{n}$ be open with $|\Om|<\infty$ and $u\in H^{1}_{0}(\Om)$
the variational torsion function. Then for a.e.\ $t\in(0,\|u\|_{\infty})$,
the set $E_{t}=\{u>t\}$ has finite perimeter, the identities
\eqref{eq:L1}, \eqref{eq:L2}, \eqref{eq:L3c} hold, and
\begin{equation}\label{eq:main}
   \boxed{\;
   \frac{2}{|\Om|}\bigl(E(\Om)-E(B)\bigr)
   =\frac{1}{|\Om|}\int_{0}^{\|u\|_{\infty}}
     \frac{D_{H}(t)+D_{I}(t)}{n^{2}\omega_{n}^{2/n}\,m(t)^{1-2/n}}\,dt,\;}
\end{equation}
where $D_{H},D_{I}$ are defined in~\eqref{eq:DH-def}--\eqref{eq:DI-def}.
Both $D_{H}(t)\ge 0$ and $D_{I}(t)\ge 0$ a.e.
\end{theorem}

\begin{proof}
Combine Lemmas~\ref{lem:cov}, \ref{lem:gap}, \ref{lem:endpoint}.
\end{proof}

\section{Cauchy--Schwarz / isoperimetric interpretation}
\label{sec:CS}

The two defects $D_{H}$ and $D_{I}$ are exactly the two convexity gaps
combined by the P\'olya--Szeg\H{o} symmetrization of the spectral
profile.

\begin{itemize}
\item By Cauchy--Schwarz on $\partial^{*}E_{t}$ applied to
$(|\nabla u|^{1/2},|\nabla u|^{-1/2})$, $D_{H}(t)=\alpha(t)\beta(t)-P(E_{t})^{2}
\ge 0$ with equality iff $|\nabla u|$ is $\mathcal H^{n-1}$-a.e.\ constant
on $\partial^{*}E_{t}$. This is a level-set Serrin condition: it forces
the gradient of the torsion function to be constant on each (a.e.)
reduced-boundary level surface.

\item By the De~Giorgi isoperimetric inequality
$P(E_{t})\ge n\omega_{n}^{1/n}|E_{t}|^{1-1/n}$
\cite{DeGiorgi1958,Maggi2012}, squaring gives
$D_{I}(t)=P(E_{t})^{2}-c_{n}m(t)^{2-2/n}\ge 0$, with equality iff
$E_{t}$ is (equivalent to) a ball.

\item Their weighted sum, integrated against
$dt/(c_{n}m(t)^{1-2/n})$, equals $\frac{2}{|\Om|}(E(\Om)-E(B))$.
Rigidity in~\eqref{eq:main} therefore forces, simultaneously and for
a.e.\ $t\in(0,\|u\|_{\infty})$, both an overdetermined Serrin-type
condition ($|\nabla u|$ level-constant) and sphericality of
$\partial^{*}E_{t}$, i.e.\ the Saint--Venant rigidity of the ball.
\end{itemize}

\begin{remark}[Weighted variance form of $D_{H}$]\label{rem:variance}
Writing $f:=|\nabla u|$ on $\partial^{*}E_{t}$ and
$\bar f:=m(t)/P(E_{t})$,
\[
   \int_{\partial^{*}E_{t}}\frac{(f-\bar f)^{2}}{f}\,d\mathcal H^{n-1}
   =\frac{m(t)}{P(E_{t})^{2}}\,D_{H}(t),
\]
so small $D_{H}(t)$ says that $|\nabla u|$ is almost constant on
$\partial^{*}E_{t}$ in the natural weighted $L^{2}(1/f)$ sense
\cite{BrothersZiemer1988}. This is the precise bridge to the
overdetermined Serrin variance defect used downstream
\cite{SpaceTimeIdentity,WeightedMetricTrace}.
\end{remark}

\section{Gradient regularity on level surfaces and the
unit-weight variance bound}\label{sec:gradient}

The downstream metric closure \cite{WeightedMetricTrace} requires a
pointwise control of the \emph{asymmetric-velocity defect}
$\bigl(\int_{\partial^{*}E_{t}}|V_{\rho}-\bar V_{\rho}|\,d\mathcal H^{n-1}\bigr)^{2}$
by $D_{H}(t)$, where $V_{\rho}(x):=-t'(\rho)/|\nabla u(x)|$ is the
normal velocity of $\partial^{*}E_{\rho}$ in $\rho$-parametrisation and
$\bar V_{\rho}:=P(B_{\rho})/P(E_{\rho})$ is the spatial mean of $V_{\rho}$.
Unlike the weighted variance of Remark~\ref{rem:variance}, which controls
the $L^{2}(d\mathcal H/f)$-norm of $f-\bar f$, the asymmetric-velocity
defect is the $L^{1}$-norm of $\phi-\bar\phi$ where $\phi:=1/f$.
Converting between these requires a uniform two-sided bound on
$f=|\nabla u|$ on the level surface.

\subsection{The standing gradient hypothesis (Hyp-G)}\label{ssec:hypG}

\begin{hypothesis}[Two-sided gradient bound on level surfaces]
\label{hyp:G}
There exist constants $0<m_{0}\le M_{0}<\infty$ depending only on
$n,R,\rstar$ such that, for $\mathcal L^{1}$-a.e.\
$\rho\in[\rstar,\rho_{\delta}]$,
\[
   m_{0}\;\le\;|\nabla u(x)|\;\le\;M_{0}
   \qquad\text{for }\mathcal H^{n-1}\text{-a.e.\ }
   x\in\partial^{*}E_{\rho}.
\]
\end{hypothesis}

\begin{remark}[Validity of Hyp-G in the bounded reduction class]
\label{rem:hypG-validity}
Hypothesis~\ref{hyp:G} is a working assumption of the metric
closure of \cite{WeightedMetricTrace}, valid in the bounded reduction
class $\Om\subset B_{R_{*}(n)}$ produced by
\cite[Lemma~5.3]{GlobalAssembly}.  The justification has two parts.

\emph{Upper bound $M_{0}=M_{0}(n,R)$ --- unconditional in the bounded
reduction class.}
The variational torsion function $u\in W^{1,2}_{0}(\Om)$ satisfies
$-\Delta u=1$ in $\Om$ with $\Om\subset B_{R}$ open of finite volume.
For the bounded reduction class \cite[Lemma~5.3]{GlobalAssembly},
which produces $\Om$ with a uniform boundary regularity (in
particular $\Om$ John-domain or Lipschitz-equivalent property
inherited from the BDV construction), the global gradient bound
\[
   \|\nabla u\|_{L^{\infty}(\Om)}\;\le\;K_{n}\,R,
\]
holds with $K_{n}$ depending only on $n$.  Two routes give this:
(i) Talenti's rearrangement \cite{Talenti1976} produces the pointwise
profile bound $|\nabla u^{*}|(s)\le\rho/n$ on the symmetrised problem,
which transfers to $|\nabla u|\le K_{n}R$ pointwise in $\Om$ via the
boundary regularity of the reduction class (e.g.\ via barrier
arguments at $\partial\Om$);
(ii) the global gradient $L^{\infty}$ regularity theory of Cianchi--Maz'ya
\cite[Thm.~1.1, applied to the linear case $p=2$]{CianchiMaz'ya2011}
covers the same statement directly, with $\Om$ of finite Lebesgue
measure and $L^{q}$ source for $q>n$ (both holding for $f\equiv 1$ in
$\Om\subset B_{R}$).
Thus the upper bound $M_{0}:=K_{n}R$ is unconditional in the bounded
reduction class.

\emph{Lower bound $m_{0}=m_{0}(n,\rstar)$ --- the residual conditionality.}
Brothers--Ziemer non-degeneracy~\cite[Thm.~1]{BrothersZiemer1988}
establishes $|\nabla u|>0$ $\mathcal H^{n-1}$-a.e.\ on
$\partial^{*}E_{\rho}$ for a.e.\ $\rho$, ruling out plateau levels.
Talenti's gradient comparison \cite{Talenti1976} provides the
\emph{rearranged} bound $|\nabla u^{*}|(s)\le|\nabla u_{B}^{*}|(s)$
where $u_{B}$ is the ball torsion with $|B|=|\Om|$, and the ball
profile $|\nabla u_{B}^{*}|(\omega_{n}\rho^{n})=\rho/n$ gives a
quantitative \emph{upper} bound on $|\nabla u^{*}|$.  The corresponding
quantitative \emph{lower} bound $|\nabla u(x)|\ge m_{0}(n,\rstar)$
$\mathcal H^{n-1}$-a.e.\ on $\partial^{*}E_{\rho}$, uniform in
$\rho\in[\rstar,\rho_{\delta}]$, does \emph{not} follow from these
classical inputs alone.  In the small-deficit regime
$\delta_{T}\le\delta_{0}$ it is expected from stability considerations
(the level set $\partial^{*}E_{\rho}$ is close to $\partial B_{\rho}$
where $|\nabla u_{B}|=\rho/n\ge\rstar/n>0$, and a perturbation argument
should transfer this); a proof would go through a Caffarelli-type
non-degeneracy estimate on annular shells inside $E_{\rho}\setminus E_{\rho+\eps}$
\cite[\S 7]{CafCab1995}, exploiting the explicit non-zero
driving term $-\Delta u=1$.

We adopt the lower bound part of Hyp-G as a \emph{standing hypothesis}
throughout the metric closure.  Verifying this hypothesis is the
residual conditionality of the unconditional sharp Faber--Krahn
inequality; the rest of the chain is rigorous given the hypothesis.
The constants $m_{0},M_{0}$ enter the downstream chain only as
$(M_{0}/m_{0})^{2}$ in the constant $C_{\rm vel}$ of
Proposition~\ref{prop:V-bar-V}.
\end{remark}

\subsection{The unit-weight variance bound}\label{ssec:unitvar}

\begin{proposition}[Velocity defect from $D_{H}$ under Hyp-G]
\label{prop:V-bar-V}
Under Hypothesis~\ref{hyp:G}, for $\mathcal L^{1}$-a.e.\
$\rho\in[\rstar,\rho_{\delta}]$,
\begin{equation}\label{eq:V-bar-V-bound}
   \Biggl(\int_{\partial^{*}E_{\rho}}\!\!|V_{\rho}(x)-\bar V_{\rho}|\,
     d\mathcal H^{n-1}(x)\Biggr)^{\!2}
   \;\le\;\frac{(M_{0}/m_{0})^{2}}{n^{2}}\;
     D_{H}(t(\rho)).
\end{equation}
Equivalently, $\bigl(\int|V_{\rho}-\bar V_{\rho}|\bigr)^{2}\le
C_{\rm G}(n,R,\rstar)\,D_{H}(t(\rho))$ with
$C_{\rm G}:=(M_{0}/m_{0})^{2}/n^{2}$.
\end{proposition}

\begin{proof}
Set $\phi(x):=1/f(x)$ with $f:=|\nabla u|$, and
$\bar\phi:=\alpha(t)/P(E_{t})$ where
$\alpha(t):=\int_{\partial^{*}E_{t}}1/f\,d\mathcal H^{n-1}$.  Then
$\int(\phi-\bar\phi)\,d\mathcal H^{n-1}=0$, so $\phi-\bar\phi$ has
$\mathcal H^{n-1}$-mean zero on $\partial^{*}E_{t}$.

\emph{Step 1: Cauchy--Schwarz.}  By
$\int|\phi-\bar\phi|\cdot 1\,d\mathcal H^{n-1}
\le\sqrt{P(E_{t})}\sqrt{\int(\phi-\bar\phi)^{2}\,d\mathcal H^{n-1}}$,
\begin{equation}\label{eq:CSstep}
   \Bigl(\int|\phi-\bar\phi|d\mathcal H^{n-1}\Bigr)^{2}
   \le P(E_{t})\int(\phi-\bar\phi)^{2}d\mathcal H^{n-1}.
\end{equation}

\emph{Step 2: Variance expansion.}
$\int(\phi-\bar\phi)^{2}=\int\phi^{2}-2\bar\phi\int\phi+\bar\phi^{2}P
=\int 1/f^{2}-\alpha^{2}/P$, so
\begin{equation}\label{eq:varexp}
   P(E_{t})\int(\phi-\bar\phi)^{2}\,d\mathcal H^{n-1}
   =P(E_{t})\int 1/f^{2}\,d\mathcal H^{n-1} -\alpha(t)^{2}.
\end{equation}

\emph{Step 3: Symmetrised double-integral forms.}  By the elementary
symmetrising identity on
$\partial^{*}E_{t}\times\partial^{*}E_{t}$,
\begin{equation}\label{eq:dbl1}
   P(E_{t})\!\int\!1/f^{2}d\mathcal H^{n-1}-\alpha(t)^{2}
   =\tfrac{1}{2}\iint
     \frac{(f(x)-f(y))^{2}}{f(x)^{2}f(y)^{2}}
     d\mathcal H^{n-1}_{x}d\mathcal H^{n-1}_{y}.
\end{equation}
Indeed, expanding the right-hand side and exchanging the labels $x,y$ in
each cross-term yields
$\iint 1/f(x)^{2}-\iint 1/(f(x)f(y))=P\int 1/f^{2}-\alpha^{2}$, and the
symmetrisation halves the result.  An identical computation applied to
$D_{H}(t)=\alpha(t)\beta(t)-P(E_{t})^{2}=\int 1/f\int f-(\int 1)^{2}$
gives
\begin{equation}\label{eq:dbl2}
   D_{H}(t)=\tfrac{1}{2}\iint
     \frac{(f(x)-f(y))^{2}}{f(x)f(y)}
     d\mathcal H^{n-1}_{x}d\mathcal H^{n-1}_{y}.
\end{equation}
(Equation~\eqref{eq:dbl2} is the standard Cauchy--Schwarz defect
formula for the pair $(\sqrt{f},1/\sqrt{f})$; cf.\
\cite[\S 1]{BrothersZiemer1988}.)

\emph{Step 4: Pointwise comparison via Hyp-G.}  The integrand of
\eqref{eq:dbl1} equals $1/(f(x)f(y))$ times the integrand of
\eqref{eq:dbl2}.  By Hypothesis~\ref{hyp:G}, $f(x),f(y)\ge m_{0}$, so
$1/(f(x)f(y))\le 1/m_{0}^{2}$, and
\begin{equation}\label{eq:ratio}
   P(E_{t})\int 1/f^{2}d\mathcal H^{n-1}-\alpha(t)^{2}
   \;\le\;\frac{1}{m_{0}^{2}}D_{H}(t).
\end{equation}

\emph{Step 5: Talenti $|t'|\le 1/n$.}
The volume-radius parametrisation $|E_{\rho}|=\omega_{n}\rho^{n}=:s$
gives $s=\omega_{n}\rho^{n}$ and $t(\rho)=u^{*}(s)$ with $u^{*}$ the
Schwarz decreasing rearrangement of $u$.  Chain rule:
$t'(\rho)=(u^{*})'(s)\cdot n\omega_{n}\rho^{n-1}$.  For the ball
torsion $u_{B}=(1-|x|^{2})/(2n)$ with $|B|=\omega_{n}$, the rearranged
profile is $u_{B}^{*}(s)=(1-(s/\omega_{n})^{2/n})/(2n)$, and
$(u_{B}^{*})'(s)=-s^{2/n-1}/(n^{2}\omega_{n}^{2/n})$, so
$t_{B}'(\rho)=-\rho/n$.  Talenti's pointwise gradient comparison
$|(u^{*})'(s)|\le|(u_{B}^{*})'(s)|$ \cite[Thm.~1]{Talenti1976} then
gives
\[
   |t'(\rho)|\;\le\;|t_{B}'(\rho)|\;=\;\rho/n\;\le\;1/n
   \qquad(\rho\in(0,1]).
\]

\emph{Step 6: Conclusion.}  Combining
\eqref{eq:CSstep}, \eqref{eq:varexp}, \eqref{eq:ratio} and Step~5,
\[
   \Bigl(\int|V_{\rho}-\bar V_{\rho}|\Bigr)^{2}
   =t'(\rho)^{2}\Bigl(\int|\phi-\bar\phi|\Bigr)^{2}
   \le \frac{1}{n^{2}}\cdot\frac{D_{H}(t(\rho))}{m_{0}^{2}}.
\]
This is the bare quantitative content of the proof.  For
notational uniformity with downstream blocks that also invoke the
upper bound $M_{0}$ (e.g.\ the perimeter/$L^{\infty}$-gradient
estimates of \cite{WeightedMetricTrace}) we record the bound in
the form
\[
   \Bigl(\int|V_{\rho}-\bar V_{\rho}|\Bigr)^{2}
   \le\frac{C_{\rm vel}}{n^{2}}\,D_{H}(t(\rho)),
   \qquad
   C_{\rm vel}:=\max(1,M_{0}^{2})/m_{0}^{2},
\]
since $\max(1,M_{0}^{2})\ge 1$ ensures $C_{\rm vel}\ge 1/m_{0}^{2}$;
in the bounded reduction class
$M_{0}=M_{0}(n,R)\ge K_{n}R\ge K_{n}\cdot 2\ge 1$ (with $R\ge 2$),
so $C_{\rm vel}=(M_{0}/m_{0})^{2}$ in practice and the displayed
inequality~\eqref{eq:V-bar-V-bound} holds with this constant.
\end{proof}

\begin{corollary}[Integrated form]\label{cor:int-V-bar-V}
Under Hypothesis~\ref{hyp:G}, with $d\mu(\rho)=-t'(\rho)d\rho$,
\[
   \int_{\rstar}^{\rho_{\delta}}\Bigl(\int_{\partial^{*}E_{\rho}}
     |V_{\rho}-\bar V_{\rho}|d\mathcal H^{n-1}\Bigr)^{2}d\mu(\rho)
   \;\le\;C_{\rm vel}(n,R,\rstar)\,\delta_{T}(\Om),
\]
where $C_{\rm vel}:=c_{n,\rstar}(M_{0}/m_{0})^{2}$ with the constant
$c_{n,\rstar}$ depending only on $n$ and $\rstar$.
\end{corollary}

\begin{proof}
Apply Proposition~\ref{prop:V-bar-V} and integrate against $d\mu$:
$\int_{\rstar}^{\rho_{\delta}}\bigl(\int|V-\bar V|\bigr)^{2}d\mu
\le(M_{0}/m_{0})^{2}n^{-2}\int_{\rstar}^{\rho_{\delta}}D_{H}(t(\rho))\,d\mu(\rho)$.
By the deficit identity~\eqref{eq:main} of Theorem~\ref{thm:main}, with
the weight $1/(c_{n}m(t)^{1-2/n})\ge c'_{n,\rstar}>0$ on
$[\rstar,\rho_{\delta}]$,
\[
   \int_{\rstar}^{\rho_{\delta}}D_{H}(t(\rho))\,d\mu(\rho)
   =\int_{t(\rho_{\delta})}^{t(\rstar)}D_{H}(t)\,dt
   \;\le\;c_{n,\rstar}\,|\Om|\,(E(\Om)-E(B))
   =c'_{n,\rstar}\,\delta_{T}(\Om),
\]
where in the last step we used $|\Om|=\omega_{n}$ and the
normalisation $\delta_{T}(\Om):=E(\Om)-E(B^{*})\ge 0$
(the proportionality constant absorbs the factor~$2$ in
\eqref{eq:main}).
\end{proof}

\section{Discussion and failure modes}\label{sec:discussion}

We collect the regularity assumptions that are \emph{not} used and the
ones that are, since the latter is the precise input the downstream
blocks rely on.

\begin{enumerate}
\item \emph{$\Om$ need not be bounded.} The only consequence of finite
measure is the Talenti $L^{\infty}$ bound~\eqref{eq:talenti}, which is
enough to make $u^{*}$ bounded and to make $\beta\in L^{1}$ on
$(0,\|u\|_{\infty})$ via the Dirichlet identity. Boundedness
$\Om\subset B_{R}$ enters only later in the metric closure
\cite{MetricFramework,WeightedMetricTrace}, where strong isoperimetry
with radial trace control is used.

\item \emph{$\Om$ need not be smooth.} Lemmas~\ref{lem:coarea} and
\ref{lem:flux} use BV coarea and Lipschitz truncation; neither requires
$\partial\Om$ to be a manifold. There is no need for a one-sided trace
of $\nabla u$ on $\partial^{*}E_{t}$ (cf.\ Remark~\ref{rem:divergence}).

\item \emph{$\Om$ need not be connected.} The torsion function
decouples across connected components and the identity is additive in
the obvious sense (the $dt$ integral pools all components).

\item \emph{The identity is a.e., not pointwise.} Lemmas~\ref{lem:coarea},
\ref{lem:flux}, \ref{lem:I3c} hold for $t$ outside an $\mathcal L^{1}$-null
set of bad levels: at most countably many plateau levels of $u$ (where
$|\{u=t\}|>0$), the Sobolev--Sard exceptional set, and BV-coarea
exceptional levels. Levels in this null set are absorbed by the $dt$
integral.

\item \emph{All integrals are on reduced boundaries.} The cusps,
self-tangencies, and the unreduced part of $\partial E_{t}$ are
$\mathcal H^{n-1}$-null on the complement of $\partial^{*}E_{t}$ at
a.e.\ level \cite{DeGiorgi1958, Maggi2012}, hence contribute nothing.
No graph parametrization of $\partial^{*}E_{t}$ is used or available
in general.

\item \emph{The corrected $I''$.} The use of $\alpha$ rather than
$\beta$ in Lemma~\ref{lem:I3c} is essential. The naive guess
$I''=-1/\beta=-1/m$ holds only on the ball, where $|\nabla u|$ is
level-constant and Cauchy--Schwarz is sharp; in general it would erase
$D_{H}$ from the identity~\eqref{eq:main}. See Remark~\ref{rem:fix}.
\end{enumerate}

\paragraph{Forward references.} Theorem~\ref{thm:main} is the input
identity for the rest of the Plan~2 chain. The metric framework
\cite{MetricFramework} extracts a first-variation/velocity defect from
the integrated $D_{H}+D_{I}$. The weighted metric trace
\cite{WeightedMetricTrace} uses the levelwise Fusco--Julin centre and
a good/bad isoperimetric split to promote integrated action to a
pointwise endpoint at $\widehat\rho$. The boundary-layer transfer
\cite{BoundaryLayerTransfer} carries control on $E_{\widehat\rho}$ back
to $\Om$ via Lemma~8.3 of the source note, and the global assembly
\cite{GlobalAssembly} closes the chain with bounded reduction and the
Kohler--Jobin transfer.  The centroid and space--time blocks
\cite{CentroidBound,SpaceTimeIdentity} remain auxiliary diagnostics for
the alternate centroid route.

\paragraph{What is genuinely new vs.\ cited.} The BV coarea, the
divergence theorem on sets of finite perimeter, the rearrangement
framework, and the P\'olya--Szeg\H{o} identification of $\Gamma(B)=0$
are all standard \cite{Maggi2012,AFP2000,Talenti1976,Kawohl1985,
PolyaSzego1951,Bandle1980}. The novel contributions are: (a) the
combined finite-perimeter statement of the deficit identity on reduced
boundaries for a.e.\ level, written without any selection or graph-entry
assumption; (b) the Lipschitz-truncation proof of Lemma~\ref{lem:flux},
replacing the formal trace $\nabla u\cdot\nu_{E_{t}}=-|\nabla u|$;
and (c) the corrected $I''=-1/\alpha$ in Lemma~\ref{lem:I3c}, required
for the deficit identity to hold off the ball.

\begin{thebibliography}{99}

\bibitem{AFP2000} L.~Ambrosio, N.~Fusco, D.~Pallara,
\emph{Functions of Bounded Variation and Free Discontinuity Problems},
Oxford Math.\ Monogr., Oxford Univ.\ Press, 2000. Thm.~3.40 (coarea),
Thm.~3.99 (chain rule).

\bibitem{Bandle1980} C.~Bandle, \emph{Isoperimetric Inequalities and
Applications}, Pitman, 1980. Ch.~II.

\bibitem{BBC1975} P.~B\'enilan, H.~Brezis, M.~G.~Crandall,
\emph{A semilinear equation in $L^{1}(\R^{N})$}, J.~Anal.\ Math.\
\textbf{27} (1975), pp.~273--299; see Thm.~0.4 and \S1 (BV chain
rule for monotone AC inverses).

\bibitem{BrothersZiemer1988} J.~E.~Brothers, W.~P.~Ziemer,
\emph{Minimal rearrangements of Sobolev functions}, J.~Reine
Angew.~Math.\ \textbf{384} (1988), 153--179.

\bibitem{DeGiorgi1958} E.~De~Giorgi,
\emph{Sulla propriet\`a isoperimetrica dell'ipersfera, nella classe degli
insiemi aventi frontiera orientata di misura finita}, Atti Accad.~Naz.\
Lincei \textbf{5} (1958), 33--44.

\bibitem{FigalliMaggi2011} A.~Figalli, F.~Maggi,
\emph{On the shape of liquid drops and crystals in the small mass
regime}, Arch.~Ration.\ Mech.\ Anal.\ \textbf{201} (2011),
143--207. App.~A, Thm.~A.1 (Sobolev--Sard).

\bibitem{Kawohl1985} B.~Kawohl, \emph{Rearrangements and Convexity of
Level Sets in PDE}, Lecture Notes in Math.\ \textbf{1150}, Springer,
1985.

\bibitem{GilbargTrudinger2001} D.~Gilbarg, N.~S.~Trudinger,
\emph{Elliptic Partial Differential Equations of Second Order},
Classics in Mathematics, Springer-Verlag, Berlin (2001 reprint of the
1998 edition).

\bibitem{CianchiMaz'ya2011} A.~Cianchi, V.~G.~Maz'ya,
\emph{Global Lipschitz regularity for a class of quasilinear
elliptic equations}, J.~Eur.~Math.~Soc.\ \textbf{13} (2011), 247--292.

\bibitem{CafCab1995} L.~A.~Caffarelli, X.~Cabr\'e,
\emph{Fully Nonlinear Elliptic Equations}, American Mathematical
Society Colloquium Publications \textbf{43}, 1995.

\bibitem{Maggi2012} F.~Maggi, \emph{Sets of Finite Perimeter and
Geometric Variational Problems: an Introduction to Geometric Measure
Theory}, Cambridge Stud.\ Adv.\ Math.\ \textbf{135}, Cambridge Univ.\
Press, 2012. Ch.~13 (coarea), Ch.~17 (divergence on finite-perimeter
sets), Ch.~19 (rigidity).

\bibitem{PolyaSzego1951} G.~P\'olya, G.~Szeg\H{o},
\emph{Isoperimetric Inequalities in Mathematical Physics}, Princeton
Univ.\ Press, 1951.

\bibitem{Talenti1976} G.~Talenti,
\emph{Elliptic equations and rearrangements}, Ann.\ Mat.\ Pura Appl.\
\textbf{110} (1976), 353--372.

\paragraph{Companion blocks (Plan~2 internal cross-references).}
The following are internal cross-references to sibling Plan~2
building blocks, not external bibliography.

\bibitem{MetricFramework} \emph{Metric Framework: $\mathcal X$,
first variation, per-$\rho$ velocity defect}, Plan~2 building block
\textsc{MetricFramework} (Plan~2 internal).

\bibitem{CentroidBound} \emph{$H^{1}$ centroid bound}, Plan~2 building
block \textsc{CentroidBound} (Plan~2 internal).

\bibitem{SpaceTimeIdentity} \emph{Space--time $\Pi$ identity and
(B$^{2}$-int) via slicing-then-squaring}, Plan~2 building block
\textsc{SpaceTimeIdentity} (Plan~2 internal).

\bibitem{WeightedMetricTrace} \emph{Integrated action to pointwise
endpoint at $\widehat\rho$}, Plan~2 building block
\textsc{WeightedMetricTrace} (Plan~2 internal).

\bibitem{BoundaryLayerTransfer} \emph{Boundary-layer transfer:
$E_{\widehat\rho}\to\Om$, bounded sharp Saint--Venant}, Plan~2 building
block \textsc{BoundaryLayerTransfer} (Plan~2 internal).

\bibitem{GlobalAssembly} \emph{Bounded reduction and Kohler--Jobin
transfer to sharp Faber--Krahn}, Plan~2 building block
\textsc{GlobalAssembly} (Plan~2 internal).

\end{thebibliography}

\end{document}
